\[
\text{The quotient topology on } A \text{ induced by } p
\text{ is given by}
\]

\[
\tau_A=\{U\subseteq A:p^{-1}(U)\text{ is open in }\mathbb{R}\}.
\]

Since
\[
p^{-1}(\{a\})=(-\infty,0),
\]
\[
p^{-1}(\{b\})=(0,\infty),
\]
\[
p^{-1}(\{c\})=\{0\},
\]
and
\[
p^{-1}(\{a,b\})=\mathbb{R}\setminus\{0\},
\]
we observe that the open subsets of \(A\) are
\[
\varnothing,\quad \{a\},\quad \{b\},\quad
\{a,b\},\quad \{a,b,c\}.
\]

Hence, the quotient topology on \(A\) is
\[
\boxed{
\tau_A=
\{\varnothing,\{a\},\{b\},\{a,b\},\{a,b,c\}\}.
}
\]
\]

\begin{proof}
(a) Let \(p:X\to Y\) be a continuous map. Suppose there exists a
continuous map \(f:Y\to X\) such that
\[
p\circ f=I_Y,
\]
where \(I_Y\) is the identity map on \(Y\).

We know that \(p\) is continuous. Also, since
\[
p(f(y))=y,\qquad \forall y\in Y,
\]
the map \(p\) is surjective.

Now, let \(U\subseteq Y\) be such that \(p^{-1}(U)\) is open
in \(X\).

Since \(f\) is continuous, we have
\[
f^{-1}\big(p^{-1}(U)\big)
\]
open in \(Y\).

But
\[
f^{-1}\big(p^{-1}(U)\big)
=(p\circ f)^{-1}(U)
=I_Y^{-1}(U)
=U.
\]

Therefore, \(U\) is open in \(Y\).

Thus, \(p\) is a continuous surjective map satisfying
\[
U\text{ is open in }Y
\iff p^{-1}(U)\text{ is open in }X.
\]

Hence, \(p\) is a quotient map.

\[
\boxed{p\text{ is a quotient map.}}
\]

\medskip

(b) Let \(A\subseteq X\) and let \(r:X\to A\) be a
retraction such that
\[
r(a)=a,\qquad \forall a\in A.
\]

We know that \(r\) is continuous by definition.

Also, since \(r(a)=a\) for every \(a\in A\), the map \(r\)
is surjective.

Now, let \(U\subseteq A\) be such that \(r^{-1}(U)\)
is open in \(X\).

Since \(A\) has the subspace topology, we have
\[
r^{-1}(U)\cap A
\]
open in \(A\).

But, using the property of retraction,
\[
\begin{aligned}
r^{-1}(U)\cap A
&=\{x\in A:r(x)\in U\}\\
&=\{x\in A:x\in U\}\\
&=U.
\end{aligned}
\]

Therefore, \(U\) is open in \(A\).

Thus, \(r\) is a continuous surjective map such that
\[
U\text{ is open in }A
\iff r^{-1}(U)\text{ is open in }X.
\]

Hence, \(r\) is a quotient map.

\[
\boxed{\text{Every retraction is a quotient map.}}
\]
\end{proof}

\begin{proof}
Let
\[
\pi_1:\mathbb{R}\times\mathbb{R}\to\mathbb{R}
\]
be the projection onto the first coordinate, defined by
\[
\pi_1(x,y)=x.
\]

Let
\[
A=\{(x,y)\in\mathbb{R}\times\mathbb{R}:x\geq 0
\text{ or }y=0\}.
\]

Define \(g:A\to\mathbb{R}\) by restricting \(\pi_1\) to \(A\).
Thus,
\[
g(x,y)=x.
\]

We shall show that \(g\) is a quotient map but is neither
open nor closed.

\medskip

\textbf{(i) \(g\) is a quotient map:}

Since \(g\) is the restriction of the continuous projection
map \(\pi_1\), it is continuous.

Also, \(g\) is surjective, because for every \(x\in\mathbb{R}\),
\[
(x,0)\in A
\quad\text{and}\quad
g(x,0)=x.
\]

Now, consider the map
\[
f:\mathbb{R}\to A
\]
defined by
\[
f(x)=(x,0).
\]

Clearly, \(f\) is continuous and
\[
(g\circ f)(x)=g(x,0)=x.
\]

Therefore,
\[
g\circ f=I_{\mathbb{R}}.
\]

By Question 2(a), since \(g\) and \(f\) are continuous
and \(g\circ f=I_{\mathbb{R}}\), it follows that
\[
\boxed{g\text{ is a quotient map.}}
\]

\medskip

\textbf{(ii) \(g\) is not an open map:}

Consider the set
\[
U=A\cap\big((-1,1)\times(1/2,3/2)\big).
\]

Since \((-1,1)\times(1/2,3/2)\) is open in
\(\mathbb{R}^2\), \(U\) is open in \(A\).

Now, by the definition of \(A\), if \((x,y)\in U\),
then \(x\geq 0\), since \(y\neq 0\).

Consequently,
\[
g(U)=[0,1).
\]

But \([0,1)\) is not open in \(\mathbb{R}\).

Thus, \(g\) does not map every open subset of \(A\)
onto an open subset of \(\mathbb{R}\).

Hence,
\[
\boxed{g\text{ is not an open map.}}
\]

\medskip

\textbf{(iii) \(g\) is not a closed map:}

Consider the subset
\[
F=\left\{\left(\frac{1}{n},n\right):
n\in\mathbb{N}\right\}.
\]

Clearly, \(F\subseteq A\), since \(1/n>0\) for every
\(n\in\mathbb{N}\).

We claim that \(F\) is closed in \(A\).

Indeed, \(F\) is closed in \(\mathbb{R}^2\), because
the second coordinates of its points tend to infinity,
so it has no accumulation point in \(\mathbb{R}^2\).
Therefore, \(F\) is also closed in the subspace \(A\).

Now,
\[
\begin{aligned}
g(F)
&=g\left(\left\{\left(\frac{1}{n},n\right):
n\in\mathbb{N}\right\}\right)\\
&=\left\{\frac{1}{n}:n\in\mathbb{N}\right\}.
\end{aligned}
\]

But
\[
\frac{1}{n}\longrightarrow 0
\quad\text{as }n\longrightarrow\infty,
\]
and \(0\notin g(F)\).

Hence, \(g(F)\) is not closed in \(\mathbb{R}\).

Therefore, \(g\) does not map every closed subset of
\(A\) onto a closed subset of \(\mathbb{R}\).

Thus,
\[
\boxed{g\text{ is not a closed map.}}
\]

Combining (i), (ii) and (iii), we conclude that
\[
\boxed{
g:A\to\mathbb{R}\text{ is a quotient map which is
neither open nor closed.}
}
\]
\end{proof}

\begin{proof}
\textbf{(a)} Define an equivalence relation on
\(X=\mathbb{R}^2\) by
\[
(x_0,y_0)\sim(x_1,y_1)
\iff x_0+y_0^2=x_1+y_1^2.
\]

Let \(X^*=X/\sim\) be the corresponding quotient space.

We shall show that \(X^*\) is homeomorphic to \(\mathbb{R}\).

Define a map
\[
g:X\to\mathbb{R}
\]
by
\[
g(x,y)=x+y^2.
\]

Since \(g\) is a polynomial function, it is continuous.
Also, \(g\) is surjective, since for every
\(t\in\mathbb{R}\),
\[
g(t,0)=t.
\]

Now, define a map
\[
f:\mathbb{R}\to X
\]
by
\[
f(t)=(t,0).
\]

Clearly, \(f\) is continuous and
\[
(g\circ f)(t)=g(t,0)=t.
\]

Thus,
\[
g\circ f=I_{\mathbb{R}}.
\]

Therefore, by Question 2(a), \(g\) is a quotient map.

Let
\[
q:X\to X^*
\]
be the natural quotient map defined by
\[
q(x,y)=[(x,y)].
\]

Define
\[
\overline{g}:X^*\to\mathbb{R}
\]
by
\[
\overline{g}([(x,y)])=x+y^2.
\]

We show that \(\overline{g}\) is a homeomorphism.

\textit{(i) \(\overline{g}\) is well-defined:}

Suppose
\[
[(x_0,y_0)]=[(x_1,y_1)].
\]

Then, by the definition of the equivalence relation,
\[
x_0+y_0^2=x_1+y_1^2.
\]

Hence,
\[
\overline{g}([(x_0,y_0)])
=
\overline{g}([(x_1,y_1)]).
\]

Therefore, \(\overline{g}\) is well-defined.

\textit{(ii) \(\overline{g}\) is bijective:}

Suppose
\[
\overline{g}([(x_0,y_0)])
=
\overline{g}([(x_1,y_1)]).
\]

Then
\[
x_0+y_0^2=x_1+y_1^2.
\]

Thus,
\[
(x_0,y_0)\sim(x_1,y_1),
\]
which implies
\[
[(x_0,y_0)]=[(x_1,y_1)].
\]

Therefore, \(\overline{g}\) is injective.

Further, since \(g\) is surjective, for every
\(t\in\mathbb{R}\), there exists \((x,y)\in X\)
such that
\[
\overline{g}([(x,y)])=t.
\]

Hence, \(\overline{g}\) is surjective.

Therefore, \(\overline{g}\) is bijective.

\textit{(iii) \(\overline{g}\) is a homeomorphism:}

We have
\[
g=\overline{g}\circ q.
\]

Since \(g\) is continuous and \(q\) is a quotient map,
the quotient property of \(q\) implies that
\(\overline{g}\) is continuous.

Also, let \(U\) be any open subset of \(X^*\).
Since \(\overline{g}\) is injective, we have
\[
g^{-1}(\overline{g}(U))=q^{-1}(U).
\]

Since \(q\) is a quotient map and \(U\) is open in
\(X^*\), \(q^{-1}(U)\) is open in \(X\).

Thus, \(g^{-1}(\overline{g}(U))\) is open in \(X\).
Since \(g\) is a quotient map, \(\overline{g}(U)\)
is open in \(\mathbb{R}\).

Therefore, \(\overline{g}\) is an open map.

Hence, \(\overline{g}\) is a continuous bijective open
map, and consequently it is a homeomorphism.

Thus,
\[
\boxed{X^*\cong\mathbb{R}.}
\]

\medskip

\textbf{(b)} Define an equivalence relation on
\(X=\mathbb{R}^2\) by
\[
(x_0,y_0)\sim(x_1,y_1)
\iff x_0^2+y_0^2=x_1^2+y_1^2.
\]

Let \(X^*=X/\sim\) be the corresponding quotient space.

We shall show that \(X^*\) is homeomorphic to
\([0,\infty)\).

Define a map
\[
g:X\to[0,\infty)
\]
by
\[
g(x,y)=x^2+y^2.
\]

Since \(g\) is a polynomial function, it is continuous.
Also, \(g\) is surjective, since for every
\(t\in[0,\infty)\),
\[
g(\sqrt{t},0)=t.
\]

Now, define a map
\[
f:[0,\infty)\to X
\]
by
\[
f(t)=(\sqrt{t},0).
\]

Clearly, \(f\) is continuous and
\[
(g\circ f)(t)=g(\sqrt{t},0)=t.
\]

Thus,
\[
g\circ f=I_{[0,\infty)}.
\]

Therefore, by Question 2(a), \(g\) is a quotient map.

Let
\[
q:X\to X^*
\]
be the natural quotient map defined by
\[
q(x,y)=[(x,y)].
\]

Define
\[
\overline{g}:X^*\to[0,\infty)
\]
by
\[
\overline{g}([(x,y)])=x^2+y^2.
\]

We show that \(\overline{g}\) is a homeomorphism.

\textit{(i) \(\overline{g}\) is well-defined:}

Suppose
\[
[(x_0,y_0)]=[(x_1,y_1)].
\]

Then, by the definition of the equivalence relation,
\[
x_0^2+y_0^2=x_1^2+y_1^2.
\]

Hence,
\[
\overline{g}([(x_0,y_0)])
=
\overline{g}([(x_1,y_1)]).
\]

Therefore, \(\overline{g}\) is well-defined.

\textit{(ii) \(\overline{g}\) is bijective:}

Suppose
\[
\overline{g}([(x_0,y_0)])
=
\overline{g}([(x_1,y_1)]).
\]

Then
\[
x_0^2+y_0^2=x_1^2+y_1^2.
\]

Thus,
\[
(x_0,y_0)\sim(x_1,y_1),
\]
which implies
\[
[(x_0,y_0)]=[(x_1,y_1)].
\]

Therefore, \(\overline{g}\) is injective.

Further, since \(g\) is surjective, for every
\(t\in[0,\infty)\), there exists \((x,y)\in X\)
such that
\[
\overline{g}([(x,y)])=t.
\]

Hence, \(\overline{g}\) is surjective.

Therefore, \(\overline{g}\) is bijective.

\textit{(iii) \(\overline{g}\) is a homeomorphism:}

We have
\[
g=\overline{g}\circ q.
\]

Since \(g\) is continuous and \(q\) is a quotient map,
the quotient property of \(q\) implies that
\(\overline{g}\) is continuous.

Also, let \(U\) be any open subset of \(X^*\).
Since \(\overline{g}\) is injective, we have
\[
g^{-1}(\overline{g}(U))=q^{-1}(U).
\]

Since \(q\) is a quotient map and \(U\) is open in
\(X^*\), \(q^{-1}(U)\) is open in \(X\).

Thus, \(g^{-1}(\overline{g}(U))\) is open in \(X\).
Since \(g\) is a quotient map, \(\overline{g}(U)\)
is open in \([0,\infty)\).

Therefore, \(\overline{g}\) is an open map.

Hence, \(\overline{g}\) is a continuous bijective open
map, and consequently it is a homeomorphism.

Thus,
\[
\boxed{X^*\cong[0,\infty).}
\]

\end{proof}